LIPAS (Lunar Impact and Predictive Analysis System) is a hybrid forecasting and decision-support platform that combines machine learning with physics-based computational models to predict the Moon's space-weather conditions. These conditions include temperature, radiation, dust activity, micrometeorite impacts, solar storms, and moonquakes. By unifying space-weather data products from NASA and the National Oceanic and Atmospheric Administration (NOAA) alongside OMNI2 data, the LIPAS platform processes the data through a temporal Long Short-Term Memory (LSTM) neural network. Users can view the LIPAS space-weather predictions via an interactive, web-based dashboard for any lunar location.

This study addresses three research questions. First, it documents how the hybrid LIPAS model was developed and what was learned in that process. Second, it examines whether a unified multi-hazard system can support operational decision-making in a controlled mission simulation used alongside existing public NASA and NOAA products. Third, it tests whether the locked hybrid improves prediction for selected hazards relative to selected replicated public baselines and physics-only references. The quantitative evaluation uses a summer Window A (1 June--4 August 2026) and a disturbed Window B (20 April--20 May 2024) under the locked serve blend on the second training-run LSTM weights (blend [0.40, 0.15, 0.02, 0.02, 0.03, 0.01]; radiation share 0.15). Hybrid RMSE is lower than the replica on all six hazards in both windows; micrometeorite and dust contrasts are tempered for circular / near-circular label risk. The contribution evaluated here is a hazard-specific hybrid architecture inside a lunar decision-support framework; LIPAS is assessed as a forecasting tool used alongside specialist NASA and NOAA models, not as a replacement for them.
